%%%%%%%%%%%%%%% CARACTÈRES SPÉCIAUX %%%%%%%%%%%%%%%

%% V/, R/, A/, la croix, l'astérisque et l'obèle, pour usage dans GABC (<sp>V/</sp>, <sp>R/</sp>, <sp>A/</sp>, <sp>cross</sp>, <sp>*</sp>, <sp>+</sp>)
\newfontfamily\plantinsymbols[
  Path = ./,
  Extension = .ttf,
  UprightFont = *-Regular,
  ItalicFont  = *-Italic,
  Ligatures = TeX,
]{PlantinSymbols}
\gresetspecial{V/}{\textcolor{gregoriocolor}{\plantinsymbols{V.~}}}
\gresetspecial{R/}{\textcolor{gregoriocolor}{\plantinsymbols{R.~}}}
\gresetspecial{+}{\textcolor{gregoriocolor}{\fontspec{Plantin MT Pro}†~}}
\gresetspecial{*}{\textcolor{gregoriocolor}{\gresixstar}}
\gresetspecial{++}{\textcolor{gregoriocolor}{\grecross}}
\gresetspecial{crux}{\textcolor{gregoriocolor}{\fontspec{FreeSerif}\symbol{"2720}}}
\gresetspecial{lcrux}{\textcolor{gregoriocolor}{+}}
%% Les mêmes, pour usage dans LaTeX (\vv \rr \aa \cc \psstar \pscross), pour usage dans le texte et les psaumes
\newcommand{\specialcharhsep}{3mm} % space after invoking R/ or V/ or A/ outside rubrics
\newcommand{\vv}{\textcolor{gregoriocolor}{\plantinsymbols{V.}\nolinebreak[4]\hspace{\specialcharhsep}\nolinebreak[4]}}
\newcommand{\rr}{\textcolor{gregoriocolor}{\plantinsymbols{R.}\nolinebreak[4]\hspace{\specialcharhsep}\nolinebreak[4]}}
\newcommand{\cc}{\textcolor{gregoriocolor}{\fontspec{FreeSerif}\symbol{"2720}~}}
\newcommand{\psstar}{~\gresixstar}
\newcommand{\pscross}{~\fontspec{Plantin MT Pro}†~\allowbreak}
\newcommand{\psnonrepetitur}{\textcolor{gregoriocolor}{\fontspec{Plantin MT Pro}†~}}
%% Les mêmes, pour usage dans les rubriques (pas d'espace avant, couleur identique au reste de la rubrique)
\newcommand{\vvrub}{{\plantinsymbols{V.~}}}
\newcommand{\rrrub}{{\plantinsymbols{R.~}}}
%% Les mêmes, pour usage en annotation
\newcommand{\vvann}{{\plantinsymbols{V.}}}
\newcommand{\rrann}{{\plantinsymbols{R.}}}
%% numéros de strophes d'hymnes
\gresetspecial{2.}{\textcolor{gregoriocolor}{2.}}
\gresetspecial{3.}{\textcolor{gregoriocolor}{3.}}
\gresetspecial{4.}{\textcolor{gregoriocolor}{4.}}
\gresetspecial{5.}{\textcolor{gregoriocolor}{5.}}
\gresetspecial{6.}{\textcolor{gregoriocolor}{6.}}
\gresetspecial{7.}{\textcolor{gregoriocolor}{7.}}
\gresetspecial{8.}{\textcolor{gregoriocolor}{8.}}
\gresetspecial{Ps}{\textcolor{gregoriocolor}{Ps.}}
\gresetspecial{Cant}{\textcolor{gregoriocolor}{Cant.}}
\gresetspecial{TP}{\textcolor{gregoriocolor}{T. P.}}

%%%%%%%%%%%%%%% INTERCALAIRES

%% ligne horizontale sur toute la largeur de la page
\newcommand{\lineamagna}{
	\par
	\xleaders\vbox to 4ex{\vfill\rule{\textwidth}{0.5pt}\vfill}\vskip4ex
}

%% ligne horizontale sous les titres
\newcommand{\lineaparva}{
	\par
	\xleaders\vbox to 4ex{\vfill\hfil\rule{0.3\textwidth}{0.5pt}\hfil\vfill}\vskip4ex
}

%%%%%%%%%%%%%%% GRAPHIES SPÉCIFIQUES

%% Versicules
\newcommand{\versiculus}[2]{
	\vv #1 \\ \rr #2
}

%% Rubriques
\newcommand{\rubrica}[1]{\textcolor{gregoriocolor}{{\footnotesize #1}}}

%% Rubriques employées comme titre d'une pièce ("per annum" centré)
\newcommand{\titulumrubrica}[1]{\needspace{5\baselineskip}\begin{center}\textsc{#1}\end{center}}

%% Petit titre intermédiaire
\newcommand{\titulum}[1]{\needspace{3\baselineskip}\begin{center}\textbf{#1}\end{center}}

%% Oraison
\newcommand{\oratio}[1]{%
	%% #1 : texte de l'oraison, avec \psstar et \pscross pour les astérisques et obèles
	\needspace{3\baselineskip} % Pour assurer que le titre "Oratio" n'est pas seul en bas de page
	{\centering\textcolor{gregoriocolor}{Oratio}\par}
	\rdrop #1
}

%% Capitule
\newcommand{\capitulum}[2]{%
	%% #1 : référence biblique
	%% #2 : texte du capitule, avec \psstar et \pscross pour les astérisques et obèles
	\needspace{3\baselineskip} % Pour assurer que le titre "Capitulum" n'est pas seul en bas de page
	\hspace*{0.04\textwidth}%
	\textcolor{gregoriocolor}{Capitulum}%
	\hfill%
	\textcolor{gregoriocolor}{#1}%
	\hspace*{0.04\textwidth}%

	\rdrop #2
}

%%%%%%%%%%%%%%% COMMON RUBRICS %%%%%%%%%%%%%%%

\newcommand{\dominexaudiversiculus}{%
	\versiculus{Dómine, exáudi oratiónem meam.}{Et clamor meus ad te véniat.}%
}

\newcommand{\dominusvobiscumversiculus}{%
	\versiculus{Dóminus vobíscum.}{Et cum spíritu tuo.}%
}